% Generated by Listas/make_statements.py from pratica-lecture-04.md, item 12. Do not edit.
\textbf{Source.} Tadelis (2013), \emph{Game Theory: An Introduction}, Exercise 8.16 ``The Value of Commitment'', p.~174 {[}PDF 193{]}. Reproduced verbatim. \emph{Note (not in the book):} ``We solved the optimal consumption plan'' refers to pp.~168--169, whose printed \(x_1 = \frac{2}{3}K\) is a slip; with the objective as printed the sophisticated plan is \((\frac{1}{2}, \frac{1}{3}, \frac{1}{6})K\) (\texttt{Notas/L03/L03-08-time-inconsistency.md} §2.2). Use the corrected plan as the no-commitment benchmark in (b). Answer in the source: none in the book (the preface, p.~xiii, says about half are solved on Princeton UP's website; not checked).



\begin{quote}

\textbf{The Value of Commitment:} Consider the three-period example of a player with hyperbolic discounting described in Section 8.3.4 with utility \(\ln(x)\) in each of the three periods and with discount factors \(\delta = 1\) and \(\beta = \frac{1}{2}\). We solved the optimal consumption plan of a sophisticated player 1.



\begin{enumerate}

\def\labelenumi{\alph{enumi}.}

\item

  Imagine that an external entity can enforce any plan of action that player 1 chooses in \(t = 1\) and will prevent player 2 from modifying it. What is the plan that player 1 would choose to enforce?

\item

  Assume that \(K = 90\). Up to how much of his initial budget \(K\) will player 1 be willing to pay the external entity in order to enforce the plan you found in (a)?

\end{enumerate}

\end{quote}



\textbf{Optional extras.}

- Tadelis 8.1 and 8.2 (p.~170): mixed-strategy SPE of the voluntary Battle of the Sexes (read Fig. 8.2, p.~154; the slides' Example 2) and of mutually assured destruction (read §8.3.3, pp.~163--166, Fig. 8.11).

- Tadelis 8.9 (entry deterrence with a fixed cost, p.~172) and 8.14 (hyperbolic discounting with general \(\delta\), \(\beta\), p.~174): harder versions of items 1 and 9.

- Muñoz-Garcia Ch. 4 Ex. 2 (entry-predation, p.~156), Ex. 5--6 (sequential Pareto coordination and chicken, pp.~166, 169), Ex. 13--15 (Stackelberg variants, pp.~189, 192, 195): all with worked answers.



\begin{center}\rule{0.5\linewidth}{0.5pt}\end{center}

